%% shortname: Walrasian Equilibrium, the Welfare Theorems and the Robinson Crusoe Economy
\documentclass[11pt]{article}
\usepackage{microstyle}

\begin{document}

\lessonheader{21}{Chapter 15 (continued): equilibrium in the Edgeworth box, Pareto efficiency
and the welfare theorems; the one-consumer, one-firm (Robinson Crusoe) economy}{}


\section{Equilibrium in the Edgeworth box}

Last time we considered a pure exchange economy (no production) and depicted the economy
in the Edgeworth box. Each agent's budget line passes through the endowment $\omega$ with
slope $-p_1/p_2$.

\topic{An arbitrary price vector is not an equilibrium}
At a given price vector, agent 1 demands $x_1$ and agent 2 demands $x_2$, both on the
common budget line. If these points differ, the markets do not clear.

\begin{center}
\begin{tikzpicture}[scale=1.05]
  % box 6 x 4; agent 1: u = sqrt(x y); agent 2: (6-x)^0.3 (4-y)^0.7; omega = (4.8,1); p1/p2 = 0.3
  \draw[thick] (0,0) rectangle (6,4);
  \node[ecnode,below left] at (0,0) {$O_1$};
  \node[ecnode,above right] at (6,4) {$O_2$};
  \draw[budget] (0,2.44) -- (6,0.64);
  \draw[budget,domain=1.24:6,samples=60] plot (\x,{4.962/\x});
  \draw[budget,domain=0:5.02,samples=60] plot (\x,{4-3.954*pow(6-\x,-0.428571)});
  \node[bundle] at (4.8,1) {};    \node[ecnode,above right] at (4.8,1) {$\omega$};
  \node[bundle] at (4.067,1.22) {}; \node[ecnode,above right] at (4.067,1.22) {$x_1$};
  \node[bundle] at (2.64,1.648) {}; \node[ecnode,below left] at (2.64,1.648) {$x_2$};
  \draw[dashed] (2.64,0) -- (2.64,1.648);
  \draw[dashed] (4.067,0) -- (4.067,1.22);
  \draw[dashed] (2.64,1.648) -- (6,1.648);
  \draw[dashed] (4.067,1.22) -- (6,1.22);
  \draw[{Stealth}-{Stealth}] (2.64,-0.3) -- (4.067,-0.3);
  \node[ecnode,below] at (3.35,-0.3) {excess demand for good 1};
  \draw[{Stealth}-{Stealth}] (6.3,1.22) -- (6.3,1.648);
  \node[ecnode,right] at (6.3,1.43) {excess supply of good 2};
  \node[ecnode] at (1.3,0.5) {\scriptsize $\nearrow$ better for 1};
  \node[ecnode] at (4.9,3.5) {\scriptsize better for 2 $\swarrow$};
\end{tikzpicture}
\end{center}

Agent 1's demand for good 1 plus agent 2's exceeds $\bar\omega_1$: there is excess demand
for good 1 and excess supply of good 2 $\Rightarrow$ no equilibrium at these prices.

To figure out what the equilibrium prices are we must account for preferences
$\Rightarrow$ add preferences (indifference curves) to the box.

\topic{Equilibrium: demand $=$ supply}
In equilibrium, demand $=$ supply: the two agents' demands are the \emph{same point} of the
box, where both indifference curves are tangent to the common budget line.

\begin{center}
\begin{tikzpicture}[scale=1.05]
  % equilibrium price p1/p2 = 0.4321, x* = (3.557, 1.537)
  \draw[thick] (0,0) rectangle (6,4);
  \node[ecnode,below left] at (0,0) {$O_1$};
  \node[ecnode,above right] at (6,4) {$O_2$};
  \fill[gray!18,domain=1.367:6,samples=60] plot (\x,{5.467/\x}) -- (6,4) -- (1.367,4) -- cycle;
  \draw[budget] (0,3.074) -- (6,0.481);
  \draw[budget,domain=1.367:6,samples=60] plot (\x,{5.467/\x});
  \draw[budget,domain=0:5.21,samples=60] plot (\x,{4-3.612*pow(6-\x,-0.428571)});
  \node[bundle] at (4.8,1) {};       \node[ecnode,below right] at (4.8,1) {$\omega$};
  \node[bundle] at (3.557,1.537) {}; \node[ecnode,below left] at (3.557,1.537) {$x_1^*=x_2^*$};
  \node[ecnode] at (3.9,3.3) {better for 1};
  \node[ecnode,anchor=west] at (6.2,0.3) {slope $-p_1^*/p_2^*$};
  \draw[thin] (6.2,0.3) -- (5.6,0.65);
\end{tikzpicture}
\end{center}

\topic{Corner equilibria}
If the equilibrium is at the boundary (corner) of the box, the tangency condition need not
be satisfied. What matters is that $x_1$ and $x_2$ are on top of each other in the box, i.e.\
demand $=$ supply --- at least when there is no disposal of goods in equilibrium.

\begin{center}
\begin{tikzpicture}[scale=0.9]
  % agent 1 consumes none of good 2 ... equilibrium on the top edge
  \draw[thick] (0,0) rectangle (6,4);
  \node[ecnode,below left] at (0,0) {$O_1$};
  \node[ecnode,above right] at (6,4) {$O_2$};
  \draw[budget] (1.8,4) -- (6,0);
  \node[bundle] at (1.8,4) {};
  \node[ecnode,above] at (1.8,4) {$x_1^*=x_2^*$};
  \draw[budget,domain=1.8:5.5,samples=60] plot (\x,{4-0.952*(\x-1.8)+0.12*(\x-1.8)^2});
  \draw[budget,domain=1.8:2.85,samples=40] plot (\x,{4-3*(\x-1.8)-0.6*(\x-1.8)^2});
  \node[bundle] at (4.6,1.333) {}; \node[ecnode,below left] at (4.6,1.333) {$\omega$};
  \node[ecnode,anchor=west] at (6.2,3.0) {agent 1's indifference curve};
  \draw[thin] (6.2,3.0) -- (5.0,2.2);
  \node[ecnode,anchor=west] at (6.2,0.6) {agent 2's indifference curve};
  \draw[thin] (6.2,0.6) -- (2.6,1.2);
\end{tikzpicture}
\end{center}
At $x^*$ agent 2 consumes none of good 2 and their indifference curve is steeper than the
budget line, so it is not tangent: agent 2 would like to sell even more of good 2, but
cannot consume a negative amount. Both agents still optimize and markets clear.

\topic{Definition (Walrasian equilibrium)}
An equilibrium consists of endogenous prices $p^*$ and an allocation $(x_1^*,x_2^*)$ such
that each $x_i^*$ is agent $i$'s optimal choice at $p^*$ given wealth $p^*\cdot\omega_i$,
and markets clear: $x_{k1}^*+x_{k2}^*=\bar\omega_k$ for $k=1,2$.


\section{Pareto efficiency and the first welfare theorem}

\topic{Pareto efficiency}
An allocation is \emph{Pareto efficient} if it is impossible to make one agent better off
without hurting someone else.

Pareto efficiency is a \emph{minimal} welfare criterion: if an allocation is not Pareto
efficient, then it cannot be socially optimal, because we can make someone better off
without hurting anyone else.

\topic{The equilibrium is Pareto efficient}
In the equilibrium picture above, the set of allocations agent 1 prefers to $x^*$ (shaded)
lies strictly above the budget line, and the set agent 2 prefers lies strictly below it.
The two sets do not overlap $\Rightarrow$ it is impossible to make one agent better off
without hurting the other.

\begin{center}
\fbox{\parbox{0.86\linewidth}{\textbf{First welfare theorem.} An equilibrium allocation
$(x_1^*,x_2^*)$ is Pareto efficient (MWG: if preferences are locally non-satiated).}}
\end{center}


\section{Towards the second welfare theorem: more terminology}

\begin{flatlist}
  \item \textbf{The Pareto set:} the set of all Pareto efficient allocations (interior
        points where the two indifference curves are tangent, running from $O_1$ to $O_2$).
  \item \textbf{The set of Pareto improvements over a specific endowment $\omega$:}
        allocations where at least one agent is better off and none is hurt (the lens
        between the two indifference curves through $\omega$).
  \item \textbf{The contract curve:} the intersection of the two sets --- allocations
        that are Pareto efficient \emph{and} Pareto improvements over $\omega$.
\end{flatlist}

\begin{center}
\begin{tikzpicture}[scale=1.05]
  % Pareto set y = 6x/(21-2x); lens through omega = (4.8,1) between x = 2.216 and 4.8
  \draw[thick] (0,0) rectangle (6,4);
  \node[ecnode,below left] at (0,0) {$O_1$};
  \node[ecnode,above right] at (6,4) {$O_2$};
  \fill[gray!20,domain=2.216:4.8,samples=60]
     plot (\x,{4-3.2436*pow(6-\x,-0.428571)}) --
     plot[domain=4.8:2.216] (\x,{4.8/\x}) -- cycle;
  \draw[budget,domain=1.2:6,samples=60] plot (\x,{4.8/\x});
  \draw[budget,domain=0:5.38,samples=60] plot (\x,{4-3.2436*pow(6-\x,-0.428571)});
  \draw[line width=0.6pt,dashed,domain=0:6,samples=60] plot (\x,{6*\x/(21-2*\x)});
  \draw[adm,domain=3.376:3.787,samples=20] plot (\x,{6*\x/(21-2*\x)});
  \node[bundle] at (4.8,1) {}; \node[ecnode,below right] at (4.8,1) {$\omega$};
  % labels
  \node[ecnode,anchor=west] at (6.3,3.2) {Pareto set};
  \draw[thin] (6.3,3.2) -- (5.5,3.35);
  \node[ecnode,anchor=west] at (6.3,2.0) {Pareto improvements over $\omega$};
  \draw[thin] (6.3,2.0) -- (3.0,1.75);
  \node[ecnode,anchor=south] at (3.1,4.3) {contract curve};
  \draw[thin,-{Stealth}] (3.1,4.3) -- (3.55,1.65);
\end{tikzpicture}
\end{center}

\topic{Interpretation}
We might expect to end up at an allocation on the contract curve if there are no prices and
people simply barter with each other.

In a market economy there is no need to go around bartering with everyone. Equilibrium
prices contain all the information you need, and this facilitates a kind of coordination
between market participants.


\section{The second welfare theorem}

The second welfare theorem asks whether it is possible to achieve \emph{any} point in the
Pareto set as an equilibrium, through a redistribution of endowments (or of wealth).

The answer is \emph{yes} when preferences are convex (MWG adds continuity and
monotonicity/local non-satiation).

\begin{center}
\begin{tikzpicture}[scale=1.05]
  % target x0 = (2.5, 0.9375) on Pareto set; supporting price slope -0.375; omega' = (4, 0.375)
  \draw[thick] (0,0) rectangle (6,4);
  \node[ecnode,below left] at (0,0) {$O_1$};
  \node[ecnode,above right] at (6,4) {$O_2$};
  \draw[line width=0.6pt,dashed,domain=0:6,samples=60] plot (\x,{6*\x/(21-2*\x)});
  \draw[budget] (0,1.875) -- (5,0);
  \draw[budget,domain=0.62:6,samples=60] plot (\x,{2.34375/\x});
  \draw[budget,domain=0:4.12,samples=60] plot (\x,{4-5.239*pow(6-\x,-0.428571)});
  \node[bundle] at (2.5,0.9375) {}; \node[ecnode,above left] at (2.5,0.9375) {$x^\circ$};
  \node[bundle] at (4.8,1) {};     \node[ecnode,above right] at (4.8,1) {$\omega$};
  \node[bundle] at (4,0.375) {};   \node[ecnode,below] at (4,0.3) {$\omega'$};
  \draw[-{Stealth},thick,shorten >=2pt] (4.8,1) -- (4,0.375);
  \node[ecnode,anchor=west] at (6.3,3.2) {Pareto set};
  \draw[thin] (6.3,3.2) -- (5.5,3.35);
  \node[ecnode,anchor=west] at (6.3,0.6) {redistribute $\omega\to\omega'$};
  \draw[thin] (6.3,0.6) -- (4.5,0.65);
\end{tikzpicture}
\end{center}

To reach the Pareto efficient allocation $x^\circ$: take the common tangent line of the two
indifference curves at $x^\circ$ as the price line, and redistribute endowments from
$\omega$ to some $\omega'$ on that line. At these prices and endowments, $x^\circ$ is an
equilibrium.


\section{The Robinson Crusoe economy (MWG 15.C)}

\topic{The agents and markets}
One consumer/worker and one firm.

\topic{Two commodities}
Labor/leisure and a consumption good.
\begin{flatlist}
  \item The consumer/worker has preferences that are monotone and convex over leisure
        $x_1$ and the consumption good $x_2$.
  \item The consumer/worker has a time endowment $\bar L$. If leisure is $x_1$, labor is
        $z=\bar L-x_1$.
  \item The firm uses labor to produce the consumption good as an output; the production
        function is $f(z)$, increasing and concave.
\end{flatlist}

\topic{The firm's problem}
Given output price $p>0$ and input price (wage) $w>0$,
\[ \max_{z\ge0}\ p f(z)-wz \qquad\text{or}\qquad \max_{z\ge0}\ f(z)-\frac{w}{p}\,z. \]
When the solution is interior, the first-order condition is
\[ \underbrace{f'(z)}_{\substack{\text{MRT: how we can transform}\\ \text{input into output}}}
   \;=\; \underbrace{\frac{w}{p}}_{\text{relative price}}. \]
The profit function is $\pi(p,w)=$ maximized profit.

\topic{The consumer/worker's problem}
The consumer/worker has two sources of income: earnings from labor $+$ dividend (profit)
as the owner of the firm. The budget constraint is
\begin{align*}
  p\,x_2 &\le w(\bar L-x_1)+\pi(p,w)\\
  \Longleftrightarrow\quad
  \underbrace{w\,x_1}_{\substack{\text{opportunity cost:}\\\text{loss of leisure}}}+\,p\,x_2
   &\le \underbrace{w\bar L+\pi(p,w)}_{\text{potential income}}\\
  \Longleftrightarrow\quad
  \frac{w}{p}\,x_1+x_2 &\le \underbrace{\frac{w}{p}}_{\text{real wage}}\bar L
   +\underbrace{\frac{\pi(p,w)}{p}}_{\text{real dividend}}.
\end{align*}
Now maximize utility subject to the budget constraint. When the solution is interior,
\[ \text{MRS}=\frac{w}{p}. \]

\topic{In combination}
\[ \boxed{\ \text{MRT}=\frac{w}{p}=\text{MRS}\ } \]

\begin{center}
\begin{tikzpicture}[scale=0.95]
  % f(z) = 2 sqrt(z), Lbar = 4, u = sqrt(x1 x2): x1* = 8/3, x2* = 2.309, w/p = 0.866
  \ecaxes{5.2}{5.2}
  \node[ecnode,right] at (5.2,0) {$x_1$ (leisure)};
  \node[ecnode,above] at (0,5.2) {$x_2$};
  \draw[budget,domain=0:4,samples=80] plot (\x,{2*sqrt(4-\x)});
  \draw[budget] (0,4.619) -- (4,1.155);
  \draw[dashed] (4,0) node[ecnode,below] {$\bar L$} -- (4,1.155);
  \draw[dashed] (4,1.155) -- (4.3,1.155) node[ecnode,right] {$\pi/p$};
  \draw[budget,domain=1.45:4.9,samples=60] plot (\x,{6.158/\x});
  \node[bundle] at (2.667,2.309) {};
  \draw[dashed] (2.667,0) node[ecnode,below] {$x_1^*$} -- (2.667,2.309);
  \draw[{Stealth}-{Stealth}] (2.667,-0.55) -- (4,-0.55);
  \node[ecnode,below] at (3.33,-0.55) {$z^*$};
  \node[ecnode,anchor=west] at (0.15,4.85) {slope $-w/p$};
  \node[ecnode,anchor=west] at (0.15,2.3) {$x_2=f(\bar L-x_1)$};
  \draw[thin] (0.6,2.5) -- (0.9,3.4);
  \node[ecnode,anchor=west] at (1.75,4.3) {indifference curve};
\end{tikzpicture}
\end{center}

$\Rightarrow$ The outcome is efficient in the sense that the economy's capacity to
transform one good into the other is equal to the willingness to trade one good for
another.

$\Rightarrow$ The first welfare theorem with production.

\topic{Compare to a planner}
A planner who wants to
\[ \max_{x_1,x_2}\ u(x_1,x_2)\quad\text{s.t.}\quad x_2=f(\bar L-x_1) \]
has the first-order condition
\[ \frac{\partial u/\partial x_1}{\partial u/\partial x_2}=f'(\bar L-x_1),
   \qquad\text{i.e.}\quad \text{MRS}=\text{MRT}, \]
the same condition as the market equilibrium.

\runin{Next:} MWG 17.A--17.B, 17.C.


\begin{transcriptnotes}
  \item The export has 6 pages; page 6 has only a few lines (the rest is blank). No date or
        title is written, so the header date is left empty; the subtitle and all section
        headings are \textbf{added} (MWG Chapter 15, Sections 15.B--15.C).
  \item Page 1: ``dipicted'' $\to$ \emph{depicted}; ``edgeworth box'' $\to$
        \emph{Edgeworth box}; ``excess demand for Good 1, excess supply for good 2
        $\Rightarrow$ no equl'' $\to$ \emph{no equilibrium}. The page sketch (budget line
        through $\omega$, demands $x_1$, $x_2$, arrows for the directions of improvement)
        is redrawn to scale with Cobb--Douglas preferences ($u_1=\sqrt{x_{11}x_{21}}$,
        $u_2=x_{12}^{0.3}x_{22}^{0.7}$), $\omega_1=(4.8,1)$ in a $6\times4$ box and
        $p_1/p_2=0.3$; the braces marking the excess demand/supply are \textbf{added}.
        The sentence explaining why the two demands violate market clearing is
        \textbf{added}.
  \item Page 1 lower half: ``Add preference indiff curve to the box''; ``equlibrium''
        $\to$ \emph{equilibrium}. The equilibrium sketch is redrawn with the same
        preferences at the equilibrium price $p_1^*/p_2^*\approx0.432$ (the page's
        unlabelled point is labelled $x_1^*=x_2^*$).
  \item Page 1, right column: ``effcivency'', ``minimal welfare criterion'', ``if the
        allocation is not pareto effuent, the it cannot be socilly optimal \ldots\ make
        someone better of without hurting any one'' $\to$ corrected spelling; ``the''
        $\to$ \emph{then}. The arrow from this sentence points into a shaded lens in the
        second (corner) sketch; that lens is the set of Pareto improvements, discussed
        properly on page 2, so it is not repeated in the corner figure.
  \item Page 1, corner sketch: ``In eq is at the corner then tangency condition need not be
        satisfied, $x_1$ and $x_2$ are on top of each other in the box or demand $=$
        supply'' $\to$ \emph{If the equilibrium is at the corner \ldots}. Redrawn
        schematically with the equilibrium on the top edge (agent 1 tangent, agent 2 at a
        corner with a steeper indifference curve); the explanatory sentence below the
        figure is \textbf{added}. Page 2 opens with ``At least when there's no
        disposal of goods in equilibrium'', which continues this sentence and is placed
        here.
  \item Page 2: ``Equilibrium consist of endogenous $p^*$ And $(x_1^*,x_2^*)$''; a struck
        ``de'' is dropped. The full definition (optimality and market clearing) is
        \textbf{added} as the ``Definition'' paragraph. ``EQ is pareto effiiet'' $\to$
        \emph{Pareto efficient}. The page's shaded region (agent 1's better set above the
        budget line) is shown in the equilibrium figure of Section 1; the sentence about
        the two preferred sets not overlapping is \textbf{added}.
  \item ``First welfare theorem: EQ is pareto efficient''; the condition ``locally
        non-satiated'' is \textbf{added} from MWG (Prop.\ 16.C.1).
  \item Page 2 lower half: ``teminology'' $\to$ \emph{terminology}; ``set of pareto
        improvement over A specific endowment, at least agent is better off, none are
        hurt'' $\to$ \emph{at least one agent}. Pareto-set sketch redrawn: with the
        preferences above the Pareto set is $x_{21}=6x_{11}/(21-2x_{11})$ (dashed), the
        lens is shaded and the contract curve is the thick segment
        $x_{11}\in[3.38,3.79]$. The page also draws several tangent indifference-curve
        pairs along the Pareto set; these are omitted.
  \item Page 3: ``The contract curve intersection of the two sets'';
        ``pareto improwent over $w$'' $\to$ \emph{Pareto improvements over $\omega$}.
        ``bartening with each [other]'' $\to$ \emph{barter with each other}; ``with
        [struck: equivlente] everyone'' (struck word dropped). ``this failitates''
        $\to$ \emph{facilitates}.
  \item Page 3 lower half: ``asks if it is possible to achieve any point in the pareto set
        through redistribution (of endowments) on wealth'' $\to$ \emph{or wealth};
        ``as an equilibrium'' is \textbf{added}. ``The answer is when preferences are
        convex'' $\to$ \emph{The answer is yes when \ldots}; MWG's additional conditions
        are \textbf{added} in parentheses. The sketch (target point on the Pareto set,
        arrow moving $\omega$ along the supporting price line) is redrawn to scale; the
        labels $x^\circ$, $\omega'$ and the explanation below the figure are
        \textbf{added}.
  \item Page 4: ``Rohinso crusoe'' $\to$ \emph{Robinson Crusoe}; ``montone''
        $\to$ \emph{monotone}. The time endowment is written $L$; typeset as $\bar L$ to
        avoid confusion with labor. A crossed-out symbol before ``the consumer/worker has
        time endowment'' is dropped. The line ``good as an output the production is
        $f(z)$, increasing and concave'' straddles the page halves and is reconstructed.
  \item Page 4: the second objective is written ``$\max_z \ell(z)-\frac{w}{p}z$''; the
        $\ell$ is read as $f$ (dividing the first objective by $p$). ``when solution is
        written'' $\to$ \emph{interior}. ``foc'' $\to$ \emph{first-order condition};
        ``How we can transform input'' completed to \emph{into output}.
  \item Page 5: ``have'' $\to$ \emph{has}; ``[ear]nings'' and ``[th]e firm'' are cut off at
        the left margin; ``divident'' $\to$ \emph{dividend}; ``opportun.'' $\to$
        \emph{opportunity cost}. The page writes ``$w_1x_1$''; \textbf{corrected} to
        $w\,x_1$ (there is only one wage). The third form of the budget constraint is cut
        at the page break (``$\frac{w}{p}x_1+x_2 = \frac{w}{p}L+\frac{\pi(p,w)}{p}$''
        partly visible); typeset with $\le$ to match the lines above.
  \item Page 5: ``when interiorn'' $\to$ \emph{when interior}. ``In combination
        MRT $=w/p=$ MRS'' and the boxed ``MRT $=$ MRS''. The leisure--consumption diagram
        ($f(z)=2\sqrt z$, $\bar L=4$, $u=\sqrt{x_1x_2}$) is \textbf{added} to illustrate
        the tangency; it is not on the page.
  \item Page 6: ``The first welfare theorem with production''. In a side box: ``compare to a
        planner who wants $u(x_1,x_2)$ s.t.\ $x_2=f(L-x_1)$''; the first-order condition
        and the conclusion MRS $=$ MRT are \textbf{added}. ``Next 17A--B, 17C''.
\end{transcriptnotes}

\end{document}
